\begin{table}[H]
\centering
\caption{Pruebas de raíz unitaria estacional en el IPC}
\label{tab:estacional}
\small
\begin{tabular}{@{}lrrcrrc@{}}
\toprule
 & \multicolumn{3}{c}{OCSB} & \multicolumn{3}{c}{Canova--Hansen} \\
\cmidrule(lr){2-4}\cmidrule(l){5-7}
Muestra & Estad. & VC 5\% & Rechaza $H_0$ & Estad. & VC 5\% & Rechaza $H_0$ \\
\midrule
2005--2026 & $-8.25$ & $-1.80$ & Sí & $1.76$ & $2.71$ & No \\
2009--2026 & $-7.40$ & $-1.80$ & Sí & $1.61$ & $2.71$ & No \\
\bottomrule
\end{tabular}
\notascuadro{\textit{Nota.} OCSB (Osborn et al., 1988), sobre $100 \times \log IPC$: $H_0$ = hay raíz unitaria estacional. Canova y Hansen (1995), sobre la inflación mensual: $H_0$ = la estacionalidad es determinística y estable. Las dos apuntan a estacionalidad determinística, por lo que la inflación entra en variación mensual con dummies estacionales en lugar de en variación anual (Kilian y Lütkepohl, 2017, sección 19.2). VC = valor crítico.}
\end{table}
